\chapter{Revisão Bibliográfica}

\section*{Visão Computacional Aplicada e Sensoriamento}
As aplicações práticas de visão computacional demonstram versatilidade em diferentes setores. No monitoramento biológico, Cardoso (2020) \cite{Cardoso2020} validou métodos eficazes para a contagem automatizada de espécimes em ambientes agropecuários, embora destaque a sensibilidade à iluminação como um desafio persistente. No campo da interação e acessibilidade, Cunha (2022) \cite{Cunha2022} comprovou a viabilidade de sistemas de rastreamento ocular de baixo custo utilizando hardware convencional, reforçando a tendência de democratização de soluções de VC que antes exigiam sensores de alto custo.

\section*{Robótica e Controle Assistido por Visão}
A integração entre percepção visual e controle robótico é explorada tanto em sistemas aéreos quanto terrestres e industriais. Menezes (2013) \cite{Menezes2013} contribui com a análise de controladores alternativos para servovisão em robôs aéreos, apontando melhorias de desempenho em cenários de voo. No âmbito da robótica educativa e de pesquisa, Rezeck (2019) \cite{Rezeck2019} estabeleceu bases para plataformas abertas de baixo custo, enquanto Bertoncini (2022) \cite{Bertoncini2022} e Santos (2022) \cite{Santos2022} avançaram na automação aplicada: o primeiro focando em centrais de controle para navegação autônoma e o segundo na precisão de manipuladores robóticos industriais para tarefas de pick-and-place.

\section*{Infraestrutura IoT, Segurança e Inteligência de Dados}
A sustentação dessas tecnologias depende de redes eficientes e seguras. Pereira (2022) \cite{Pereira2022} demonstra como o Aprendizado Profundo pode ser aplicado na classificação de tráfego em redes definidas por software (SDN) integradas à IoT, essencial para a Indústria 4.0. Paralelamente, a eficiência energética e a segurança surgem como pilares críticos; enquanto Silva (2023) \cite{GuimaraesDaSilva2023} oferece um comparativo rigoroso de modelos de DL para previsão de consumo de energia em dispositivos IoT, Guimarães (2023) \cite{BritoGuimaraes2023} aborda a vulnerabilidade desses sistemas, utilizando hardware de borda (Edge Computing) para a detecção de botnets em tempo real.